\chapter{Modbus TCP: Your Data, Your Handlers}
\label{ch:modbus}\index{Modbus}

The networking chapter built a socket facade: \texttt{GNAT.Sockets} over the
W5500 (Chapter~\ref{ch:networking}). This chapter puts an industrial protocol on
top of it. Modbus is the lingua franca of industrial I/O,
and Modbus TCP is the protocol over a socket: a seven-byte MBAP header (transaction id, a zero protocol id, a
length, a unit id) followed by a PDU (a function code and its data). The data
model is four tables: coils and discrete inputs (one bit each, read-write and
read-only), holding and input registers (sixteen bits each, read-write and
read-only). A slave that cannot honour a request replies with the function code's
high bit set and an \emph{exception code} (illegal function, illegal data
address, \ldots). All of which is small enough that the framing fits on one page
and is checked, byte-for-byte against the canonical example frame, by a host unit
test before any socket exists.

Two design questions shaped the API, and both were answered by the grain of this
codebase rather than by what other Modbus libraries do.

\section{The slave owns no registers}

Most Modbus libraries allocate register arrays for you and copy values in and out.
This one stores \emph{nothing}. You derive from \texttt{Modbus.Slave.Server}, keep
your data as fields of your own type, and override the handlers you support:

\begin{lstlisting}
type Device is new Modbus.Slave.Server with record
   Holding : Word_Array (0 .. 63) := (others => 0);   -- your storage
end record;

overriding procedure On_Read_Holding_Registers
  (Self : in out Device; Unit : Unit_Id; Addr : Address; Qty : Positive;
   Into : out Word_Array; Status : out Exception_Code)
is begin
   if Natural (Addr) + Qty - 1 > Self.Holding'Last then
      Status := Illegal_Data_Address;              -- you decide what's valid
   else
      for I in 0 .. Qty - 1 loop Into (Into'First + I) := Self.Holding (Natural (Addr) + I); end loop;
      Status := None;
   end if;
end On_Read_Holding_Registers;
\end{lstlisting}

This is \emph{dispatching}, not callbacks, and the choice is deliberate. With a tagged
type your storage lives in \texttt{Self}: the object \emph{is} the context, so there
is no \texttt{Ctx} pointer to thread and no closure to build, which matters under the
\texttt{No\_Implicit\_Dynamic\_Code} restriction that forbids the trampolines a nested
access-to-subprogram callback would need. Handlers you don't override default to
returning \texttt{Illegal\_Function}, so a device implements exactly the function codes
it offers and the server turns every non-\texttt{None} status into the correct
exception reply. A handler returning \texttt{Illegal\_Data\_Address} for an address it
doesn't serve is the whole of its error handling; the wire is the library's problem.

The protocol itself is factored into a socket-free \texttt{Process} (one request ADU in,
one reply ADU out, dispatching to your \texttt{Server}), so framing, dispatch, the
exception path and the \texttt{Illegal\_Function} default are all host-tested against
crafted frames with no sockets at all. \texttt{Run} is a thin one-client-at-a-time loop
around it.

\texttt{Process} treats the request bytes as \emph{untrusted}, because on a TCP port
they are. Two disciplines make a malformed frame a clean exception reply rather than
a fault. First, every field is bounds-checked before it is read: a request is
rejected with \texttt{Illegal\_Data\_Value} unless the received length actually
covers the address, quantity and (for the write-multiple functions) byte-count and
data it is about to parse --- so a truncated frame can never be interpreted against
stale bytes left in the reused receive buffer from an earlier request. Second, and
subtler, the request's 16-bit quantity is \emph{range-checked before} it is used to
size anything. The reply for a read is built in a stack array whose length is the
requested count; sizing that array straight from the wire would let a single packet
asking for 65\,535 registers allocate a 128\,KB frame and overflow the task
stack --- a one-packet remote denial of service. \texttt{Process} validates the
quantity against the per-function Modbus maximum \emph{first} and only then, in a
nested block, declares the now-bounded array. The rule generalises past Modbus: never
let an attacker-supplied count drive an allocation before you have checked it.

\section{The master returns status, not exceptions}

A master is synchronous (send a request, wait for the reply), so each call returns
its data into a caller buffer plus a \texttt{Status}:

\begin{lstlisting}
Read_Holding_Registers (S, Unit => 1, Addr => 0, Qty => 5,
                        Into => W, Result => R, Exc => E);
\end{lstlisting}

The error model follows this book's settled convention (\S\ref{ch:net:gnatsockets}):
status for \emph{expected} outcomes, Ada exceptions only for \emph{programming} errors.
A Modbus exception reply is an expected outcome: \texttt{Result} is
\texttt{Exception\_Response} and \texttt{Exc} carries the code; nothing is raised. A
timeout or a dropped connection is likewise a status. The only Ada exceptions are the
\texttt{Pre} contracts on quantity and buffer size: a caller asking for 200 registers
(the protocol caps reads at 125) or passing too small a buffer is a bug in the calling
code, caught at the call, not a runtime reply to interpret.

\section{Pinning, and a note on host-testing}

Both sides take an \emph{optional} pin. A multi-homed board (\S\ref{ch:net:multinic})
may need a slave's traffic confined to one link: the wired NIC, never cellular. The
master accepts a closure-free \texttt{Configure} hook called on the fresh socket, and
the application's hook does the pinning:

\begin{lstlisting}
procedure Pin_Eth0 (Sock : in out GNAT.Sockets.Socket_Type) is
begin GNAT.Sockets.Set_Interface (Sock, 0); end;     -- in the app, library-level
...
Modbus.Master.Connect (S, "10.0.0.7", Configure => Pin_Eth0'Access, Result => R);
\end{lstlisting}

Keeping \texttt{Set\_Interface} (a facade-only call) in the \emph{application}
rather than baked into \texttt{Modbus.Master} is what lets the master compile against a
desktop's real \texttt{GNAT.Sockets} and be host-tested end to end against a stdlib
Python slave; the slave pins by binding its listener to a specific interface's address.
The two halves were proven on a laptop (every function code, the exception path,
read-back of every write) before either touched the board, where \texttt{esp32s3\_%
modbus\_slave} (polled from a host) and \texttt{esp32s3\_modbus\_master} (pinned, polling
a host) confirmed the same against real silicon.
